%CoreSemanticsChapter.tex

%replacement text for KerML Core Semantics

%\section{8.4.3 Core Semantics}\addcontentsline{toc}{section}{8.4.3 Core Semantics}

%\section{Core Semantics Overview (KerML 8.4.3.1)}\label{8.4.3.1}
%
%\section{Core Semantic Constraints (KerML 8.4.3.1.1)}
%
%\comment{As is, for now\ldots}


\section{Core Semantics Mathematical Preliminaries (KerML 8.4.3.1.2)}\label{8.4.3.1.2}

\pn %The mathematical specification of KerML Core layer element semantics uses a model-theoretic approach. 
Core layer mathematical semantics are expressed in first order logic and set theory.

Appendix \ref{chapter:simplemath} presents an informal summary of the mathematics used to define KerML semantics.  Appendix \ref{chapter:math} presents the mathematical foundation for KerML core-layer semantics derived from the axioms of logic and set theory with utmost rigor.  Appendix \ref{chapter:KL} provides details of the modal in which the timing semantics for KerML kernel-layer are defined.  Appendix \ref{chapter:mm} lists some of the theorems and definitions from Metamath used in proofs of theorems.

%Core layer elements have no relation to time.  Consequently, their model-theoretic ontology (what is) is expressed directly in logic and set theory.
%
%Kernel layer elements' ontology uses modal logic (also expressed in logic and set theory) to provide different values for elements at different times.

%\pn For now, time will be ignored, and type, class, and feature defined as ZFC-classes.  
%In subclause (8.4.3.6)\ref{sec:classesattime} evaluating rules defining ZFC-classes at different times will be considered.

\subsection{ZFC Sets, Classes, and Well-Formed Formulas} 

\pn A \bemph{set} is a collection of values denoted by a list of elements in curly brackets:  $\{ v_1 v_2 v_3 \}$, sometimes separated by commas.

\pn A \bemph{class} (\ref{eq:df-clab}) is a collection of values denoted by a rule: $\{ \mm{x}~\vert~\mw{\varphi} \}$, which means all elements $\mm{x}$ for which rule $\mw{\varphi}$ is true.
A class will usually, but not necessarily, also be a set.
\begin{table}[h]
\caption{Set and Class Symbols}
\begin{center}
\begin{tabular}{|l|c|c|}
\hline
set construction & $\{~\dots~\}$ & set \\ \hline
class abstraction & $\{~\mm{x}~\vert~\mw{\varphi}\}$ & class \\ \hline
tuple & $\langle~\ldots~\rangle$ & tuple, sequence \\ \hline
membership & $\in$ & element (of) \\ \hline
subset(class) & $\subseteq$ & subset (of) \\ \hline
proper subset(class) & $\subset$ & subset (of) \\ \hline
empty set & $\emptyset$ & empty set \\ \hline
union & $\cup$ & union \\ \hline
intersection & $\cap$ & intersection \\ \hline
set(class) difference & $\backslash$ & difference \\ \hline
cardinality & $\mathrm{card}$ & multiplicity \\ \hline
\end{tabular}
\end{center}
\label{table:setsymbols}
\end{table}%

%\paragraph{Well-Formed Formulas}
\pn A \bemph{well-formed formula} (WFF)\index{WFF} is a boolean-valued expression conforming to a definite grammar.  
Many of the steps of proofs in Metamath prove that the theorem being proved is a WFF.  
There are proof assistants which will generate those steps automatically, or complain if what you're trying to prove isn't a WFF.
An arbitrary WFF is represented by a lower-case greek letter: $\mw{\varphi}~\mw{\psi}~\mw{\chi}~\mw{\theta}$

\pn If $\mw{\varphi}$ and $\mw{\psi}$ are WFFs, then $(\mw{\varphi}\rightarrow\mw{\psi})$ is a WFF.
\begin{table}[htp]
\caption{Logical Operators}
\begin{center}
\begin{tabular}{|l|c|c|}
\hline
%\textbf{Operation} & \textbf{Symbol} & \textbf{Text} \\ \hline
proved%\footnote{has proof; is axiom or definition; is a theorem}
 & $\vdash$ & turnstyle \\ \hline
implication & $\rightarrow$ & implies \\ \hline
biconditional & $\leftrightarrow$ & iff \\ \hline
conjunction & $\wedge$ & and \\ \hline
disjunction & $\vee$ & or \\ \hline
complement & $\lnot$ & not \\ \hline
top & $\top$ & true \\ \hline
bottom & $\bot$ & false \\ \hline
\end{tabular}
\end{center}
\label{table:logicaloperators}
\end{table}%


\subsection{Universe:  $\mathrm{V}$}

\pn The collection of all sets $\mathrm{V}$ is a class which is not a set (\ref{eq:df-v}). 
To avoid Russel's Paradox, ZFC prohibits sets having themselves as elements.  Therefore $\mathrm{V}$ which contains all sets as members, cannot include itself, so $\mathrm{V}$ cannot be a set.


\begin{restatable}{definition}{dfv}~\otherTitle{df-v}~ Define the universal class.   The letter ``V" was used earlier by Peano in 1889 for the
     universe of sets, where the letter V is derived from the word ``Verum".
     Peano's notation V was adopted by Whitehead and Russell in Principia
     Mathematica for the class of all sets in 1910.  
\begin{align}
\vdash \,\mathrm{V} = \{ \msc{x}~\vert~\msc{x} = \msc{x}
    \}\label{eq:df-v}\tag{df-v}
\end{align}
\end{restatable}

\subsection{Representation of Elements as Triples:}

\pn In the following, KerML elements will be \bemph{defined} ontologically (what is) as ZFC-classes, but they will be \bemph{represented} as a triple.

\pn Let $k$ be the kind of element (type, class, feature, datatype, class, structure, annotation, etc), $id$ be a unique identifier for the element, and $C$ be the class which defines the element:

\begin{restatable}{definition}{dfrep}~\blTitle{df-rep}~ Define KerML element representation.
\begin{align}
\langle~k,~id,~C~\rangle\label{eq:df-rep}\tag{df-rep}
\end{align}
\end{restatable}

\begin{restatable}{leandefinition}{ElementRep}~\leanTitle{ElementRep}~ Define KerML element representation in Lean.

\texttt{\textcolor{leancolor}{structure} \textcolor{leanyellow}{ElementRep} \textcolor{leancolor}{where}}

\texttt{~~k  : ElementKind}

\texttt{~~id : String}

\texttt{~~C  : Set Element}
\label{ElementRep}
\end{restatable}

\pn The kind $k$ of \ref{eq:df-rep} ranges over the fourteen KerML element kinds:

\begin{restatable}{definition}{dfkind}~\blTitle{df-kind}~ Define element kind.
\begin{align}
\vdash \mathrm{kind} = \{&~\mathrm{Element},~\mathrm{Relationship},~\mathrm{Dependency},~\mathrm{Feature},~\mathrm{Classifier},\notag\\
&~\mathrm{DataType},~\mathrm{Class},~\mathrm{Structure},~\mathrm{Association},~\mathrm{Connector},\notag\\
&~\mathrm{Behavior},~\mathrm{Function},~\mathrm{Expression},~\mathrm{Interaction}~\}\label{eq:df-kind}\tag{df-kind}
\end{align}
\end{restatable}

\begin{restatable}{leandefinition}{ElementKind}~\leanTitle{ElementKind}~ Define element kind in Lean.

\texttt{\textcolor{leancolor}{inductive} \textcolor{leanyellow}{ElementKind}}

\texttt{~~| Element | Relationship | Dependency | Feature | Classifier | DataType}

\texttt{~~| Class | Structure | Association | Connector | Behavior | Function}

\texttt{~~| Expression | Interaction}

\texttt{~~\textcolor{leancolor}{deriving} DecidableEq, Repr}
\label{ElementKind}
\end{restatable}

\subsection{Class Identifier (CI):}

\pn Triples for classes defined by \ref{eq:df-rep} have names ($id$), either a simple identifier or a qualified name which is a sequence of identifiers separated by \texttt{::}.
Regardless of whether a single identifier, or a qualified name, the whole string is considered its ``class identifier''.
The set of ZFC-class identifiers is $CI$.

\subsection{Unique Identifier (UI):} 

\pn Unique identifiers\footnote{not a UUID} will use a \bemph{wonce} which is like a \bemph{nonce} (number once), but instead is a word once:

~~~`a' through `z', then `aa' through `az', then `ba' through `bz', and so on.

\pn An instance element's unique identifier is the concatenation of its fully-qualified name made into a string (if any), ``\textbackslash\#", and a wonce.

\pn Thus even things without names will have a unique identifier:  `P::Q::\textbackslash\#abd' could be a UI for an unnamed element in namespace Q which is in namespace P.

\pn Because ``\textbackslash\#" is not among the Escape Sequences in Table 4 of KerML 8.2.2.3, there should be no conflict between unrestricted names and UI.

\pn The set of all (allocated) unique identifiers is $UI$.

\pn Let $\mathrm{mkUid}(ci)$ be a function that computes a unique identifier $ui \in UI$ from $ci \in CI$ by appending a wonce.

\pn Where $ci \in CI$ is a string (possibly empty), $tb$ is ``\textbackslash\#", $wonce$ is a wonce, and \textbf{++} is (Metamath's own, already-formalized) string concatenation:

\begin{restatable}{definition}{dfwonce}~\blTitle{df-wonce}~ Define creation of a Unique Identifier from a Class Identifier by appending a wonce.
\begin{align}
\vdash ( \,\mathrm{mkUid} ` \mc{A} ) = ( ( \mc{A} \mathbin{\operatorname{++}} \,\mathrm{tb} ) \mathbin{\operatorname{++}} \,\mathrm{wonce} )\label{eq:df-wonce}\tag{df-wonce}
\end{align}
\end{restatable}

\begin{restatable}{leandefinition}{mkUid}~\leanTitle{mkUid}~ Define creation of a Unique Identifier from a Class Identifier by appending a wonce, in Lean.

\texttt{\textcolor{leancolor}{abbrev} \textcolor{leanyellow}{CI} := String}

\texttt{\textcolor{leancolor}{abbrev} \textcolor{leanyellow}{UI} := String}

\texttt{\textcolor{leancolor}{def} \textcolor{leanyellow}{tb} : String := "\textbackslash\textbackslash\#"}

\texttt{\textcolor{leancolor}{axiom} \textcolor{leanyellow}{wonce} : String}

\texttt{\textcolor{leancolor}{noncomputable def} \textcolor{leanyellow}{mkUid} (ci : CI) : UI := ci ++ tb ++ wonce}
\label{mkUid}
\end{restatable}

%\subsection{Design}
%
%\pn Everything in a \bemph{design} $D$ is either a  type $V_T$, a  classifier $V_C$ or a feature $V_F$.
%
%\begin{definition}~\blTitle{df-design}~ Define design $D$.  
%\begin{align}
%D~=~V_T~\cup~V_C~\cup~V_F~ \notag 
%%V_T'~=~((V_T\backslash V_C)\backslash V_F)~\wedge~ \notag \\
%%V_C'~=~(V_C\backslash V_F) ~\wedge~\notag \\
%%D~=~V_T'~\cup~V_C'~\cup~V_F~\wedge \notag \\
%%V_T'\cap V_C'=\emptyset~\wedge~V_T'\cap V_F=\emptyset~\wedge~V_F\cap V_C'=\emptyset~
%\label{eq:df-design}\tag{df-design}
%\end{align}
%\end{definition}

%\chapter{Types}

\subsection{Types, Classifiers, and Features}

\pn KerML 8.4.3.1.2 defines a \bemph{vocabulary} as a triple, $V=\langle V_T, V_C, V_F\rangle$, where
\begin{description}
\item[$V_T$ ] is a set of types\index{ VT@$V_T$ types} 
\item[$V_C \subseteq V_T$ ] is a set of classifier, and\index{ VC@$V_C$ classifiers}
\item[$V_F \subseteq V_T$ ] is a set of features.\index{ VF@$V_F$ features}
\end{description}

%\pn By KerML 8.4.3.1.2, $V_T = V_C \cup V_F$, so every type must be either a classifier, or a feature.
%This doesn't seem right.  

\pn Here we assume there exist types in $V_T$ which are neither classifiers $V_C$ nor features $V_F$:
$V_C \cup V_F \subset V_T$.\footnote{KerML \S8.4.3.2 allows a \texttt{Type} that is not a \texttt{Classifier} or a \texttt{Feature}---e.g.\ an element declared with the bare \texttt{type} keyword rather than \texttt{classifier} or \texttt{feature}.  (\texttt{DataType} is \emph{not} such an example: it specializes \texttt{Classifier}, so it is already in $V_C$.)}

\pn Furthermore, we assume that classifiers and features are disjoint, $V_C \cap V_F = \emptyset$.

\subsection{Design}

\pn A KerML \bemph{design}, $D$,\index{ D@$D$ design} 
is the set of all elements in the database or workspace.  This includes all KerML libraries, standard or not.
 Design elements may be stored according to KerML 10.  In the following, concrete syntax will be used with the understanding of its correspondence to abstract syntax.


\section{Type}

\pn A \texttt{Type} is a set of values.  Therefore it is a member of the universe of sets:  $\mathrm{V}$.
The set of things classified by a type is the \bemph{extent} of the type, each
member of which is an \bemph{instance} of the type.

\pn ``\texttt{type}~$A$" represents program text where $A$ is an identifier (a String) of the
type, not the set the type itself is.  $A$ ranges over $ID$, the class of identifiers (see
\ref{eq:df-rep}).  Rather than an unstructured stand-in for ``this identifier was declared with
the \texttt{type} keyword," the concrete KerML text ``\texttt{type}~$A$" is formalized directly as
membership of a $\mathrm{Design}$-triple: $Design$ is the design (the set of all elements in the
database or workspace, \S2.1.2), and $\langle \mathsf{TYPE}, A, C \rangle \in Design$ says some
design element's triple representation (\ref{eq:df-rep}) has kind $\mathsf{TYPE}$ and identifier
$A$.  $\mathsf{TYPE}$ is deliberately \emph{not} one of \texttt{df-kind}'s fourteen enumerated
kinds (\ref{eq:df-kind}): \S2.1.1 already treats $V_C \cup V_F$ as a proper subset of $V_T$, so
$\mathsf{TYPE}$ marks a design element outside that enumeration, not a fifteenth member of it.
Given that premise, $C$ --- the type's extent, i.e.\ what $A$ denotes --- is a genuine set.

\begin{restatable}{definition}{dftype}~\blTitle{df-type}~ Define \texttt{type}.
\begin{align}
\vdash \mc{A} \in \,\mathrm{ID}\quad\&\quad \vdash \langle \mathsf{TYPE} , \mc{A} , \mc{C} \rangle
    \in \,\mathrm{Design}\quad\Rightarrow\quad \vdash \mc{C} \in
    \,\mathrm{V}\label{eq:df-type}\tag{df-type}
\end{align}
\end{restatable}

\begin{restatable}{leanaxiom}{Design}~\leanTitle{Design}~ Define \texttt{type} in Lean.
\texttt{DesignKind} is a fresh sum type, not \texttt{ElementKind}: a \texttt{Design} triple's
kind slot must hold either one of \texttt{ElementKind}'s fourteen tags or \texttt{type}, and
\texttt{type} is deliberately excluded from \texttt{ElementKind} itself.  Unlike \texttt{df-type},
no separate ``is a genuine set" assertion is needed here: \texttt{Set Element} already
guarantees \texttt{C} is a real value for every triple in \texttt{Design}, so the declarations
below are the complete translation.

\texttt{\textcolor{leancolor}{inductive} \textcolor{leanyellow}{DesignKind}}

\texttt{~~| ofElementKind (k : ElementKind) | type}

\texttt{\textcolor{leancolor}{axiom} \textcolor{leanyellow}{Design} : Set (DesignKind $\times$ Element $\times$ Set Element)}
\label{Design}
\end{restatable}

\pn Let $V_T$ be the set of all \texttt{Types} in the design (\S2.1.1's vocabulary triple $V =
\langle V_T, V_C, V_F \rangle$) --- a primitive $df\text{-}types$ \emph{constrains} rather than
fully defines, matching its own implication (not equation) form below. For every identifier $x$,
if $x$ is declared with the bare \texttt{type} keyword --- the same $Design$-triple-membership
reading \ref{eq:df-type} gives that text --- then $x \in V_T$.

\begin{restatable}{definition}{dftypes}~\blTitle{df-types}~ Define \texttt{Types}.
\begin{align}
\vdash ( \msc{x} \in \,\mathrm{ID} \rightarrow ( \exists \msc{y} \langle \mathsf{TYPE} , \msc{x} ,
    \msc{y} \rangle \in \,\mathrm{Design} \rightarrow \msc{x} \in \,\mathrm{VT}
    ) )\label{eq:df-types}\tag{df-types}
\end{align}
\end{restatable}

\begin{restatable}{leanaxiom}{VT}~\leanTitle{VT}~ Define \texttt{Types} in Lean. Unlike
\texttt{df-type} above, this genuinely needs to stay an axiom here too: \texttt{VT} is
independently primitive on both sides, so there's no automatic Lean-typing shortcut this time.

\texttt{\textcolor{leancolor}{axiom} \textcolor{leanyellow}{VT} : Set Element}

\texttt{\textcolor{leancolor}{axiom} \textcolor{leanyellow}{df\_types} : $\forall$ e : Element, ($\exists$ C : Set Element, (DesignKind.type, e, C) $\in$ Design) $\to$ e $\in$ VT}
\label{VT}
\end{restatable}

\subsection{Type Operators}

\pn A type may be defined using type operators on already defined types.  Specialization, \texttt{:>} or \texttt{specializes}, is most commonly used for subsetting, but union, intersection, difference, and disjoint are provided.  Conjugation reverses direction of features, so will be considered in \ref{sec:direction}.

\pn A specialization of a \texttt{Type} is its subset (subclass).  For all types $t_s$ and $t_g$ in $V_T$, $t_s$ specializing $t_g$ implies that $t_s$ is a subset (subclass) of $t_g$.

\begin{restatable}{definition}{dfspecializes}~\blTitle{df-specializes}~ Define \texttt{specializes}.
\begin{align}
\vdash \forall t_g , t_s \in V_T ~(\texttt{type}~t_s~\texttt{specializes}~t_g ~\rightarrow~ t_s \subseteq t_g)\label{eq:df-specializes}\tag{df-specializes}
\end{align}
\end{restatable}

\begin{restatable}{leanaxiom}{Specializes}~\leanTitle{Specializes}~ KerML's own textual \texttt{specializes} keyword,
in Lean. A fresh \texttt{Element}-level primitive: the structural elaboration of \texttt{specializes} (a
\texttt{Specialization}/\texttt{Subclassification} value) carries no semantic content of its own, and --- being a
bare, unconstrained structure --- isn't tied to what's actually asserted in the current model the way
\texttt{Design} is.

\texttt{\textcolor{leancolor}{axiom} \textcolor{leanyellow}{Specializes} : Element $\to$ Element $\to$ Prop}
\label{Specializes}
\end{restatable}

\begin{restatable}{leanaxiom}{dfSpecializes}~\leanTitle{df\_specializes}~ Define \texttt{specializes} in Lean.

\texttt{\textcolor{leancolor}{axiom} \textcolor{leanyellow}{df\_specializes} \{ts tg : Element\} \{Cs Cg : Set Element\}}

\texttt{~~(hs : (DesignKind.type, ts, Cs) $\in$ Design) (hg : (DesignKind.type, tg, Cg) $\in$ Design) :}

\texttt{~~Specializes ts tg $\to$ Cs $\subseteq$ Cg}
\label{df_specializes}
\end{restatable}

\pn The subclass relationship is defined in \ref{eq:df-ss}.

\pn When a type specializes multiple supertypes, (\ref{eq:df-specializes}) applies to \emph{each} generalization independently, so $t_s$ must be a subset of \emph{every} one of them---hence a subset of their \emph{intersection}.

\begin{restatable}{definition}{dfmultspec}~\blTitle{df-multspec}~ Define multiple \texttt{specializes}.
\begin{align}
\vdash \forall t_g ,~ \ldots,~ t_j, t_s \in V_T ~(\texttt{type}~t_s~\texttt{specializes}~t_g,~ \ldots, ~t_j ~\rightarrow~ t_s \subseteq t_g \cap \ldots \cap t_j)\label{eq:df-multspec}\tag{df-multspec}
\end{align}
\end{restatable}

\begin{restatable}{leantheorem}{dfMultspecLean}~\leanTitle{df\_multspec}~ Multiple specialization bounds the
specific type's extension within the intersection of every one of its generalizations' extensions, in Lean.
Proved directly from \texttt{df\_specializes} applied to each generalization independently, exactly as the
prose above describes it --- so, unlike \texttt{Specializes}/\texttt{df\_specializes}, this one genuinely is a
theorem, not an axiom.

\texttt{\textcolor{leancolor}{theorem} \textcolor{leanyellow}{df\_multspec} \{ts tg th : Element\} \{Cs Cg Ch : Set Element\}}

\texttt{~~(hs : (DesignKind.type, ts, Cs) $\in$ Design) (hg : (DesignKind.type, tg, Cg) $\in$ Design)}

\texttt{~~(hh : (DesignKind.type, th, Ch) $\in$ Design)}

\texttt{~~(hsg : Specializes ts tg) (hsh : Specializes ts th) : Cs $\subseteq$ Cg $\cap$ Ch :=}

\texttt{~~fun \_ hx $\Rightarrow$ $\langle$df\_specializes hs hg hsg hx, df\_specializes hs hh hsh hx$\rangle$}
\label{df_multspec}
\end{restatable}

\pn The intersection relationship is defined in \ref{eq:df-in}.


\pn The \texttt{unions} of \texttt{Types} is the set of elements in at least one of them.  Although just a pair of \texttt{Types} is used in the definition,
\texttt{unions} generalizes to any number of \texttt{Types}.  Unlike ordinary multiple specialization (\ref{eq:df-multspec}), which only bounds $t_s$ within the \emph{intersection} of its generalizations, \texttt{unions} asserts that $t_a$ \emph{equals} the union of its unioning types (KerML \S8.4.3.2, rule~3)---so \texttt{unions} is \emph{not} superfluous.

\begin{restatable}{definition}{dfunions}~\blTitle{df-unions}~ Define \texttt{unions}.
\begin{align}
\vdash \forall t_a , t_b, t_c \in V_T~(\texttt{type}~t_a~\texttt{specializes}~t_b~\texttt{unions}~t_c~ \rightarrow~ t_a~=~t_b \cup t_c)\label{eq:df-unions}\tag{df-unions}
\end{align}
\end{restatable}

\begin{restatable}{leanaxiom}{Unions}~\leanTitle{Unions}~ KerML's own textual \texttt{unions} keyword, in Lean.
A fresh \texttt{Element}-level primitive, same reasoning as \texttt{Specializes}: the structural elaboration of
\texttt{unions} (a \texttt{Unioning} value) carries no semantic content of its own and isn't tied to what's
actually asserted in the current model.

\texttt{\textcolor{leancolor}{axiom} \textcolor{leanyellow}{Unions} : Element $\to$ Element $\to$ Prop}
\label{Unions}
\end{restatable}

\begin{restatable}{leanaxiom}{dfUnionsLean}~\leanTitle{df\_unions}~ Define \texttt{unions} in Lean.

\texttt{\textcolor{leancolor}{axiom} \textcolor{leanyellow}{df\_unions} \{ta tb tc : Element\} \{Ca Cb Cc : Set Element\}}

\texttt{~~(ha : (DesignKind.type, ta, Ca) $\in$ Design) (hb : (DesignKind.type, tb, Cb) $\in$ Design)}

\texttt{~~(hc : (DesignKind.type, tc, Cc) $\in$ Design) :}

\texttt{~~Specializes ta tb $\to$ Unions ta tc $\to$ Ca = Cb $\cup$ Cc}
\label{df_unions}
\end{restatable}

%\begin{definition} Define union of two classes.
%\begin{align}
%(\mc{A}~\cup~\mc{B})~=~\{~x~\vert ~(x \in \mc{A}~\vee~x \in\mc{B})~\} \label{eq:df-un}\tag{df-un}
%\end{align}
%\end{definition}

\pn The \texttt{intersects} of \texttt{Types} is the set of elements in all of them.  
%Although just a pair of \texttt{Types} is used in the definition,
%\texttt{intersects} generalizes to any number of \texttt{Types}.

\begin{restatable}{definition}{dfintersects}~\blTitle{df-intersects}~  Define \texttt{intersects}.
\begin{align}
\vdash \forall t_a , t_b, t_c \in V_T~(\texttt{type}~t_a~\texttt{specializes}~t_b~\texttt{intersects}~t_c~ \rightarrow ~t_a~=~t_b \cap t_c)\label{eq:df-intersects}\tag{df-intersects}
\end{align}
\end{restatable}

\begin{restatable}{leanaxiom}{Intersects}~\leanTitle{Intersects}~ KerML's own textual \texttt{intersects}
keyword, in Lean. Same reasoning as \texttt{Unions} above.

\texttt{\textcolor{leancolor}{axiom} \textcolor{leanyellow}{Intersects} : Element $\to$ Element $\to$ Prop}
\label{Intersects}
\end{restatable}

\begin{restatable}{leanaxiom}{dfIntersectsLean}~\leanTitle{df\_intersects}~ Define \texttt{intersects} in Lean.

\texttt{\textcolor{leancolor}{axiom} \textcolor{leanyellow}{df\_intersects} \{ta tb tc : Element\} \{Ca Cb Cc : Set Element\}}

\texttt{~~(ha : (DesignKind.type, ta, Ca) $\in$ Design) (hb : (DesignKind.type, tb, Cb) $\in$ Design)}

\texttt{~~(hc : (DesignKind.type, tc, Cc) $\in$ Design) :}

\texttt{~~Specializes ta tb $\to$ Intersects ta tc $\to$ Ca = Cb $\cap$ Cc}
\label{df_intersects}
\end{restatable}

%The intersection relationship is defined in \ref{eq:df-in}.
%\begin{definition} Define intersection of two classes.  
%\begin{align}
%(\mc{A}~\cap~\mc{B})~=~\{~x~\vert ~(x \in \mc{A}~\wedge~x \in\mc{B})~\} \label{eq:df-in}\tag{df-in}
%\end{align}
%\end{definition}
\pn The intersection relationship is defined in \ref{eq:df-in}.

\pn The \texttt{differences} of \texttt{Types} is class difference.

\begin{restatable}{definition}{dfdifferences}~\blTitle{df-differences}~ Define \texttt{differences}.
\begin{align}
\vdash \forall t_a , t_b, t_c \in V_T~(\texttt{type}~t_a~\texttt{specializes}~t_b~\texttt{differences}~t_c~ \rightarrow ~t_a~=~(t_b~\backslash~ t_c))~\label{eq:df-differences}\tag{df-differences}
\end{align}
\end{restatable}

\begin{restatable}{leanaxiom}{Differences}~\leanTitle{Differences}~ KerML's own textual \texttt{differences}
keyword, in Lean. Same reasoning as \texttt{Unions} above.

\texttt{\textcolor{leancolor}{axiom} \textcolor{leanyellow}{Differences} : Element $\to$ Element $\to$ Prop}
\label{Differences}
\end{restatable}

\begin{restatable}{leanaxiom}{dfDifferencesLean}~\leanTitle{df\_differences}~ Define \texttt{differences} in
Lean.

\texttt{\textcolor{leancolor}{axiom} \textcolor{leanyellow}{df\_differences} \{ta tb tc : Element\} \{Ca Cb Cc : Set Element\}}

\texttt{~~(ha : (DesignKind.type, ta, Ca) $\in$ Design) (hb : (DesignKind.type, tb, Cb) $\in$ Design)}

\texttt{~~(hc : (DesignKind.type, tc, Cc) $\in$ Design) :}

\texttt{~~Specializes ta tb $\to$ Differences ta tc $\to$ Ca = Cb $\setminus$ Cc}
\label{df_differences}
\end{restatable}

\pn The class difference relationship is defined in \ref{eq:df-dif}.


\pn \texttt{Types} are \texttt{disjoint} when their intersection is the empty set (KerML \S8.4.3.2, rule~2): $t_a$ must still be a subset of $t_b$ per \texttt{specializes}, but need only have empty intersection with $t_c$---not equal a class difference.

\begin{restatable}{definition}{dfdisjoint}~\blTitle{df-disjoint}~  Define \texttt{disjoint}.
\begin{align}
\vdash \forall t_a , t_b, t_c  \in V_T~(\texttt{type}~t_a~\texttt{specializes}~t_b~\texttt{disjoint}~\texttt{from}~t_c ~\rightarrow~  ~t_a \subseteq t_b~\wedge~t_a \cap t_c = \emptyset)\label{eq:df-disjoint}\tag{df-disjoint}
\end{align}
\end{restatable}

\begin{restatable}{leanaxiom}{Disjoint}~\leanTitle{Disjoint}~ KerML's own textual \texttt{disjoint from}
keyword, in Lean. Same reasoning as \texttt{Unions} above.

\texttt{\textcolor{leancolor}{axiom} \textcolor{leanyellow}{Disjoint} : Element $\to$ Element $\to$ Prop}
\label{Disjoint}
\end{restatable}

\begin{restatable}{leanaxiom}{dfDisjointLean}~\leanTitle{df\_disjoint}~ Define \texttt{disjoint} in Lean. The
first conjunct below (\texttt{Ca} $\subseteq$ \texttt{Cb}) is, notably, already exactly \texttt{df\_specializes}'s
own conclusion given \texttt{Specializes ta tb}; it's restated here as part of the conjunction only because
that's the literal shape \texttt{df-disjoint} itself asserts above, not because it's independently new content.

\texttt{\textcolor{leancolor}{axiom} \textcolor{leanyellow}{df\_disjoint} \{ta tb tc : Element\} \{Ca Cb Cc : Set Element\}}

\texttt{~~(ha : (DesignKind.type, ta, Ca) $\in$ Design) (hb : (DesignKind.type, tb, Cb) $\in$ Design)}

\texttt{~~(hc : (DesignKind.type, tc, Cc) $\in$ Design) :}

\texttt{~~Specializes ta tb $\to$ Disjoint ta tc $\to$ Ca $\subseteq$ Cb $\wedge$ Ca $\cap$ Cc = $\emptyset$}
\label{df_disjoint}
\end{restatable}

\subsection{Type Operator Composition}

\pn \texttt{Type} abstract syntax (KerML \S8.3.3.1.1) defines \texttt{Specialization}, \texttt{Disjoining}, \texttt{Unioning}, \texttt{Intersec\-ting}, and \texttt{Differencing} as \texttt{Relationship}.
\texttt{Type} concrete syntax (KerML \S8.2.4.1.1) requires at least one \texttt{SpecializationPart} followed by any number of \texttt{TypeRelationshipPart}, which are disjoining, differencing, 
intersecting, and unioning.
\texttt{Type} abstract syntax does not appear to have any ordering among its type relationships.  

\pn Are these the same?

\begin{lstlisting}
type A specializes B unions C intersects D differences E disjoint from F
\end{lstlisting}
\[ A \subseteq (((( B \cup C) \cap D)\backslash E)\backslash F) \]

\begin{lstlisting}
type A specializes B intersects D differences E unions C disjoint from F
\end{lstlisting}
\[ A \subseteq (((( B \cap D)\backslash E)\cup C)\backslash F) \]

\begin{lstlisting}
type A specializes B differences E disjoint from F unions C intersects D
\end{lstlisting}
\[ A \subseteq (((( B\backslash E)\backslash F)  \cup C) \cap D)\]

\pn Certainly not.  Therefore the abstract syntax must preserve ordering of type operators in the concrete syntax.\footnote{ISSUE}

\subsection{Type Multiplicity}

\pn When declaring a type as a specialization of another type, the number of elements in the new type can be restricted to a \bemph{multiplicity}, a.k.a. cardinality.  

\pn $\mathrm{card} ` \mc{A}$ is application of the cardinality operator to a class in Metamath which will be abbreviated as  $\#\mc{A}$, or $\#\{x~\vert~some rule~\}$ when using class abstraction, or $\#\{a, b, c\}$ when a set is defined by listing its elements.

\begin{restatable}{definition}{dfcardinality}~\blTitle{df-cardinality}~ Define class cardinality (\#).
\begin{align}
\vdash \#\mc{A}~\equiv \mathrm{card} ` \mc{A}~\label{eq:df-cardinality}\tag{df-cardinality}
\end{align}
\end{restatable}

\begin{restatable}{leandefinition}{cardElem}~\leanTitle{cardElem}~ Define class cardinality in Lean. \texttt{\#}
is literal notation for \texttt{Set.ncard} (Mathlib's own finite-set cardinality --- "junk value zero if
infinite", matching the book's own "nothing in a system is infinite" stance), not a new primitive:
\texttt{SFS.mm}'s own \texttt{df-cardinality} is itself just an abbreviation reusing base \texttt{set.mm}'s
already-developed cardinal number theory, so this is the direct Lean analogue, not a shortcut --- a genuine
\texttt{def}, no axiom needed.

\texttt{\textcolor{leancolor}{noncomputable def} \textcolor{leanyellow}{cardElem} (A : Set Element) : $\mathbb{N}$ := A.ncard}
\label{cardElem}
\end{restatable}

\pn Multiplicity for types is defined by placing the value in square brackets, [ ].

\begin{restatable}{definition}{dftypemultiplicity}~\blTitle{df-typemultiplicity}~ Define type multiplicity
\begin{align}
\texttt{type B[n] specializes A;} \Rightarrow \notag \\
\vdash (~\mc{B}~\subseteq \mc{A}~\wedge~n\le\#\mc{A}~\wedge~\#\mc{B}=n~)
\label{eq:df-typemultiplicity}\tag{df-typemultiplicity}
\end{align}
\end{restatable}

\begin{restatable}{leanaxiom}{TypeMultiplicityExact}~\leanTitle{TypeMultiplicityExact}~ KerML's exact
multiplicity bracket \texttt{[n]}, in Lean. A fresh \texttt{Element}-level primitive, same reasoning as
\texttt{Specializes}: nothing currently represents a type's declared multiplicity.

\texttt{\textcolor{leancolor}{axiom} \textcolor{leanyellow}{TypeMultiplicityExact} : Element $\to$ $\mathbb{N}$ $\to$ Prop}
\label{TypeMultiplicityExact}
\end{restatable}

\begin{restatable}{leanaxiom}{dfTypemultiplicityLean}~\leanTitle{df\_typemultiplicity}~ Define type
multiplicity in Lean. The first conjunct below (\texttt{Cb} $\subseteq$ \texttt{Ca}) is, notably, already
exactly \texttt{df\_specializes}'s own conclusion given \texttt{Specializes B A}; it's restated here as part
of the conjunction only because that's the literal shape \texttt{df-typemultiplicity} itself asserts above,
not because it's independently new content.

\texttt{\textcolor{leancolor}{axiom} \textcolor{leanyellow}{df\_typemultiplicity} \{A B : Element\} \{Ca Cb : Set Element\} \{n : $\mathbb{N}$\}}

\texttt{~~(hA : (DesignKind.type, A, Ca) $\in$ Design) (hB : (DesignKind.type, B, Cb) $\in$ Design) :}

\texttt{~~Specializes B A $\to$ TypeMultiplicityExact B n $\to$}

\texttt{~~~~Cb $\subseteq$ Ca $\wedge$ n $\le$ cardElem Ca $\wedge$ cardElem Cb = n}
\label{df_typemultiplicity}
\end{restatable}

\pn Nothing in a system is infinite.  Therefore [*] is undefined multiplicity--not infinite multiplicity.

\pn Undefined multiplicity must be specialized into something defined, or is in error.

\pn \texttt{type B[*] specializes A;} means $\mc{B}~\subseteq \mc{A}$.

\pn Multiplicity range constrains cardinality to that range.

\begin{restatable}{definition}{dftypemultiplicityrange}~\blTitle{df-typemultiplicityrange}~ Define type multiplicity range
\begin{align}
\texttt{type B[m..n] specializes A;} \Rightarrow \notag \\
\vdash (~\mc{B}~\subseteq \mc{A}~\wedge~n\le\#\mc{A}~~\wedge~m\le\#\mc{B}~\wedge~\#\mc{B}\le n~)
\label{eq:df-typemultiplicityrange}\tag{df-typemultiplicityrange}
\end{align}
\end{restatable}

\begin{restatable}{leanaxiom}{TypeMultiplicityRange}~\leanTitle{TypeMultiplicityRange}~ KerML's multiplicity
range bracket \texttt{[m..n]}, in Lean. Same reasoning as \texttt{TypeMultiplicityExact} above.

\texttt{\textcolor{leancolor}{axiom} \textcolor{leanyellow}{TypeMultiplicityRange} : Element $\to$ $\mathbb{N}$ $\to$ $\mathbb{N}$ $\to$ Prop}
\label{TypeMultiplicityRange}
\end{restatable}

\begin{restatable}{leanaxiom}{dfTypemultiplicityrangeLean}~\leanTitle{df\_typemultiplicityrange}~ Define type
multiplicity range in Lean. Same reasoning as \texttt{df\_typemultiplicity} above.

\texttt{\textcolor{leancolor}{axiom} \textcolor{leanyellow}{df\_typemultiplicityrange} \{A B : Element\} \{Ca Cb : Set Element\} \{m n : $\mathbb{N}$\}}

\texttt{~~(hA : (DesignKind.type, A, Ca) $\in$ Design) (hB : (DesignKind.type, B, Cb) $\in$ Design) :}

\texttt{~~Specializes B A $\to$ TypeMultiplicityRange B m n $\to$}

\texttt{~~~~Cb $\subseteq$ Ca $\wedge$ n $\le$ cardElem Ca $\wedge$ m $\le$ cardElem Cb $\wedge$ cardElem Cb $\le$ n}
\label{df_typemultiplicityrange}
\end{restatable}


\section{Features Semantics (KerML 8.4.3.4)}

\pn A \texttt{Feature}\index{feature} is a relation--a set of pairs.  The set of first elements of the pairs of a \texttt{Feature} is the \bemph{domain} of the relation and \bemph{featuring type} of the \texttt{Feature}.  The set of second elements of the pairs of a \texttt{Feature} is the \bemph{co-domain} of the relation and \bemph{featured type}.

\pn Let $V_F$ be the set of all \texttt{Features} in the design $D$:  $V_F \subset D$, then for all elements in the design, an element being defined a \texttt{feature} implies the element is a member of $V_F$.

\begin{restatable}{definition}{dffeature}~\blTitle{df-feature}~  Define \texttt{feature}.
\begin{align}
\forall f \in D~(\texttt{feature}~f~ \rightarrow ~f \in V_F)~\label{eq:df-feature}\tag{df-feature}
\end{align}
\end{restatable}

\begin{restatable}{leanaxiom}{VF}~\leanTitle{VF}~ Define the set of \texttt{Features} in Lean. Exact parallel of
\texttt{VT} above, one level down.

\texttt{\textcolor{leancolor}{axiom} \textcolor{leanyellow}{VF} : Set Element}
\label{VF}
\end{restatable}

\begin{restatable}{leanaxiom}{dfFeatureLean}~\leanTitle{df\_feature}~ Define \texttt{feature} in Lean. Unlike
\texttt{type}, \texttt{Feature} already has a slot in \texttt{ElementKind}'s fourteen tags, so no new
\texttt{DesignKind} case is needed --- exact parallel of \texttt{df\_types} above.

\texttt{\textcolor{leancolor}{axiom} \textcolor{leanyellow}{df\_feature} : $\forall$ e : Element,}

\texttt{~~($\exists$ C : Set Element, (DesignKind.ofElementKind ElementKind.Feature, e, C) $\in$ Design) $\to$ e $\in$ VF}
\label{df_feature}
\end{restatable}

Note:  $V_F$ only contains stand-alone  \texttt{Features}.  Most  \texttt{Features} are declared within the \texttt{TypeBody} of a type declaration.

\subsection{Featured By}  

\pn A \texttt{Feature} can be declared as a stand-alone entity.

\pn Where Y and Z are names of \texttt{Types}:
\ecaption{Featured By}
\begin{lstlisting}
feature x typed by Y featured by Z;
\end{lstlisting}
says $x$ is a class of pairs where the first element is from $\mc{Z}$ and the second element is from $\mc{Y}$.

\ref{eq:df-featuredby} defines a relation with domain $\mc{Z}$ and co-domain $\mc{Y}$ without constraining its cardinality.


\begin{restatable}{definition}{dffeaturedby}~\blTitle{df-featuredby}~   Define \texttt{featured by}
\begin{align}
\texttt{feature}~x~\texttt{typed by}~ Y ~\texttt{featured by}~ Z~\texttt{;}~\Rightarrow  \notag \\
\vdash ~x=\{~\langle \textbf{feature},``x",\mc{C}\rangle~\vert~\mc{C}~=~\{~\langle z,y\rangle~\vert~y \in \mc{Y} \wedge z \in \mc{Z}~\}~\}~\label{eq:df-featuredby}\tag{df-featuredby}
\end{align}
\end{restatable}

\begin{restatable}{leanaxiom}{FeatureContent}~\leanTitle{FeatureContent}~ A feature's design-triple content, in
Lean, as a set of ordered pairs (a relation) --- \texttt{Design}'s own content slot (\texttt{Set Element})
can't hold that, so this is a dedicated primitive, independent of \texttt{Design}, rather than forcing this
into the \texttt{(DesignKind, Element, Set Element)} triple shape used everywhere else.

\texttt{\textcolor{leancolor}{axiom} \textcolor{leanyellow}{FeatureContent} : Element $\to$ Set (Element $\times$ Element) $\to$ Prop}
\label{FeatureContent}
\end{restatable}

\begin{restatable}{leanaxiom}{FeaturedByDecl}~\leanTitle{FeaturedByDecl}~ \texttt{feature x typed by Y featured
by Z;} in Lean. A fresh \texttt{Element}-level primitive, same reasoning as \texttt{Specializes}: nothing
currently represents this concrete-syntax clause.

\texttt{\textcolor{leancolor}{axiom} \textcolor{leanyellow}{FeaturedByDecl} : Element $\to$ Element $\to$ Element $\to$ Prop}
\label{FeaturedByDecl}
\end{restatable}

\begin{restatable}{leanaxiom}{dfFeaturedbyLean}~\leanTitle{df\_featuredby}~ Define \texttt{featured by} in Lean.

\texttt{\textcolor{leancolor}{axiom} \textcolor{leanyellow}{df\_featuredby} \{x y z : Element\} \{Cy Cz : Set Element\}}

\texttt{~~(hy : (DesignKind.type, y, Cy) $\in$ Design) (hz : (DesignKind.type, z, Cz) $\in$ Design) :}

\texttt{~~FeaturedByDecl x y z $\to$ FeatureContent x \{p : Element $\times$ Element $\vert$ p.1 $\in$ Cz $\wedge$ p.2 $\in$ Cy\}}
\label{df_featuredby}
\end{restatable}

Note:  $\mc{Z}$ doesn't ``know'' about \texttt{Feature} $x$; something else needed a relation with domain $\mc{Z}$ and co-domain $\mc{Y}$.


\subsection{Domain-less Feature}  

\pn A \texttt{Feature} can be defined without a domain, which defaults to the Universe $\mathrm{V}$.

\begin{restatable}{definition}{dfdomainless}~\blTitle{df-domainless}~   Define domain-less \texttt{Feature}
\begin{align}
\texttt{feature}~x~\texttt{typed by}~ Y ~\texttt{;}~\Rightarrow  \notag \\
\vdash~x=\{~\langle \textbf{feature},``x",\mc{C}\rangle~\vert~\mc{C}~=~\{~\langle z,y\rangle~\vert~y \in \mc{Y} \wedge z \in \mc{\mathrm{V}}~\}~\}~\label{eq:df-domainless}\tag{df-domainless}
\end{align}
\end{restatable}

\begin{restatable}{leanaxiom}{DomainlessDecl}~\leanTitle{DomainlessDecl}~ \texttt{feature x typed by Y;} (no
\texttt{featured by} clause) in Lean. A separate primitive from \texttt{FeaturedByDecl} above, mirroring the
book's own two separate definitions for the two distinct concrete-syntax forms, rather than one primitive
plus a ``no featured-by clause'' side condition.

\texttt{\textcolor{leancolor}{axiom} \textcolor{leanyellow}{DomainlessDecl} : Element $\to$ Element $\to$ Prop}
\label{DomainlessDecl}
\end{restatable}

\begin{restatable}{leanaxiom}{dfDomainlessLean}~\leanTitle{df\_domainless}~ Define domain-less \texttt{Feature}
in Lean. The universal domain quantification collapses to no constraint on the pair's first component at
all --- the same ``vacuous existential/quantifier disappears'' simplification used for
\texttt{RegionSurface}/\texttt{RegionFilm} in \S0.8.

\texttt{\textcolor{leancolor}{axiom} \textcolor{leanyellow}{df\_domainless} \{x y : Element\} \{Cy : Set Element\}}

\texttt{~~(hy : (DesignKind.type, y, Cy) $\in$ Design) :}

\texttt{~~DomainlessDecl x y $\to$ FeatureContent x \{p : Element $\times$ Element $\vert$ p.2 $\in$ Cy\}}
\label{df_domainless}
\end{restatable}

\ref{eq:df-domainless} defines a relation with domain $\mc{\mathrm{V}}$, the universe, and co-domain $\mc{Y}$ 

%\comment{Philosophical Rant:  Domain-less features are endemic in .kerml libraries.  Then, most feature declarations \texttt{subsets} some domain-less feature.
%Thus, understanding a feature declaration, requires looking in two places--at least.  The penchant for using inheritance means multiple levels of indirection.
%Indiscriminate inheritance may be entirely understandable to those with a UML mindset, but is frightful for users.  I'm certain that popularity of Rust comes from 
%prohibition of inheritance more than memory safety.  If KerML had a language construct for \texttt{self}, there would less need for feature chaining or \texttt{this} and \texttt{that}.  
%Cyclic definitions could also be \bemph{prohibited}. }
%\begin{lstlisting}
%	abstract classifier Anything {
%		doc
%		/*
%	     * Anything is the top level generalized type in the language. 
%	     */
%		
%		feature self: Anything[1] subsets things chains things.that {
%			doc
%			/*
%			 * The source of a SelfLink of this thing to itself. self is thus a feature that
%			 * relates everything to itself. It is also the value of the nested "that" feature
%			 * of all other things featured by this thing.
%			 */
%		}
%	}
%\end{lstlisting}
%\begin{lstlisting}
%	abstract feature things: Anything [1..*] nonunique {
%		doc
%		/*
%		 * things is the top-level feature in the language.
%		 */
%
%		feature that : Anything[1] {
%			doc
%			/*
%			 * For each value of things, the "featuring instance" of that value. 
%			 * This is enforced by declaring Anything::self to be the chaining of things.that, 
%			 * restricting it the single value of self.
%			 */			
%		}
%	}
%\end{lstlisting}
%\comment{THERE IS ABSOLUTELY NOTHING WHICH SAYS \texttt{that} REFERS TO THE DOMAIN OF \texttt{things} OTHER THAN DOC ANNOTATIONS!
%If \texttt{Anything} is the top-level generalized type, why is it declared as \texttt{classifier}?
%Rant off.}


\section{Types with Features (KerML 8.4.3.5)}

\pn Type specialization was defined by (\ref{eq:df-specializes}) to be subset.  However, having feature members of types makes types concrete.  
If types are sets of values in $\mathrm{V}$, what is the set defined by the following?

\begin{lstlisting}
type Z specializes W { feature x typed by Y; }
\end{lstlisting}

\pn For this, \texttt{Z} will be defined by a \bemph{rule}.  A collection defined by a rule is a \bemph{class} (in ZFC set theory [\ref{sec:sets}]).
To do so, there must be some way to refer to the identifier of the feature in the rule defining a class.  To do so, by convention we will use the type identifier as the variable used to define the class, and refer to a feature of the type by a period between the type identifier and feature identifier.


\subsection{Feature Member:}  
\pn Usually \texttt{Features} are defined as a \texttt{FeatureMember} in \texttt{TypeBody} (KerML 8.2.4.1).

\ecaption{Type with Feature Member}
\begin{lstlisting}
type Person specializes Anything { feature familyname typed by String; }
\end{lstlisting}
says type \emph{Person} has a feature named \emph{familyname} which is a relation having domain \emph{Person} and co-domain \emph{String}.

\pn First of all, given default multiplicity of one, the relation of feature \emph{familyname} will have only a single pair, $\{~\langle p,s\rangle~\}$ where $p \in Person$ and $s \in String$.
Expressing the class of all possible Person(s) requires using (\ref{eq:df-rep}).

\pn The type \emph{Person} will be the class of all triples:

\[ 
\mc{Person}=\{~\langle$\textbf{type}$,w,\langle \mc{C},\mc{F}\rangle\rangle~\vert~
\begin{array}{c}
w=W($``Person''$)~\wedge~  \mc{C}= \{~c~\vert ~c\in \mc{Anything}\}~\wedge~b \in \mc{String}~\\
\wedge~ \mc{F}~=~\{ ~\langle$\textbf{feature}$,$``familyname''$,\langle w,b\rangle\rangle~\}
\end{array}
 ~\} 
\]

\pn So, $Person$ is a class of triples, the first element is \textbf{type}, the second element is wonce of the string ``Person'', and the third element is a pair of classes $\langle C,F\rangle$.  $C$ is the class of all $c$ such that $c$ is $Anything$.
$F$ is a set containing a single triple for the \texttt{familyname} feature having domain $w$ and co-domain $b$ which holds the family name string.
%\pn Importantly, the domain of the relation $w$ is the same value as the containing triple.  This makes the relation apply to that particular person.  Here, $w$ performs the function of \texttt{self} applied to a type, and \texttt{that} applied to a feature.

\begin{restatable}{definition}{dffeaturemember}~\blTitle{df-featuremember}~   Define \texttt{Feature Member}
\begin{align}
\texttt{type}~ Z~\texttt{specializes}~T~\{ \texttt{feature}~ x ~\texttt{typed by}~ Y~\texttt{;} \}~\Rightarrow  \notag \\
\mc{Z}=\{~\langle \textbf{type},w,\langle \mc{C},\mc{F}\rangle\rangle ~\vert~
w=W(``Z")~ \wedge
\mc{C}=\{~c~\vert ~c \in \mc{T}~\}~\wedge~\notag \\
~b\in \mc{Y}~\wedge~\mc{F}~=~\{~\langle\textbf{feature},``x",\langle w,b \rangle\rangle~\}~\}~\label{eq:df-featuremember}\tag{df-featuremember}
\end{align}
\end{restatable}

\begin{restatable}{leanaxiom}{FeatureMemberOf}~\leanTitle{FeatureMemberOf}~ A type or feature $Z$'s owned
feature member, in Lean: $Z$ has a feature member identified by $x$ whose relation content is $R$. The shared
primitive \texttt{df-featuremember}/\texttt{df-featureoffeature}/\texttt{df-typebody} all build on below.

\texttt{\textcolor{leancolor}{axiom} \textcolor{leanyellow}{FeatureMemberOf} : Element $\to$ Element $\to$ Set (Element $\times$ Element) $\to$ Prop}
\label{FeatureMemberOf}
\end{restatable}

\begin{restatable}{leanaxiom}{OwnedFeatureTypedBy}~\leanTitle{OwnedFeatureTypedBy}~ \texttt{type Z \{ feature x
typed by Y; \}} (an owned feature member, no explicit \texttt{featured by}), in Lean. A fresh
\texttt{Element}-level primitive, same reasoning as \texttt{Specializes}.

\texttt{\textcolor{leancolor}{axiom} \textcolor{leanyellow}{OwnedFeatureTypedBy} : Element $\to$ Element $\to$ Element $\to$ Prop}
\label{OwnedFeatureTypedBy}
\end{restatable}

\begin{restatable}{leanaxiom}{dfFeaturememberLean}~\leanTitle{df\_featuremember}~ Define \texttt{Feature
Member} in Lean: a type $Z$ specializing $T$ with an owned feature $x$ typed by $Y$ has a feature member whose
relation pairs $T$'s extension (domain, matching the book's own $C=\{c \vert c\in T\}$ clause) with $Y$'s
extension (co-domain).

\texttt{\textcolor{leancolor}{axiom} \textcolor{leanyellow}{df\_featuremember} \{Z T x Y' : Element\} \{Ct Cy : Set Element\}}

\texttt{~~(ht : (DesignKind.type, T, Ct) $\in$ Design) (hy : (DesignKind.type, Y', Cy) $\in$ Design) :}

\texttt{~~Specializes Z T $\to$ OwnedFeatureTypedBy Z x Y' $\to$}

\texttt{~~~~FeatureMemberOf Z x \{p : Element $\times$ Element $\vert$ p.1 $\in$ Ct $\wedge$ p.2 $\in$ Cy\}}
\label{df_featuremember}
\end{restatable}

\pn This says that \texttt{type Z} is the class of all tuples where the first element is \textbf{type} (indicating it is a type), the second element is the wonce of the string ``Z'', and the third element is itself a pair of class $\langle \mc{C},\mc{F}\rangle$ where $\mc{C}$ has type $\mc{T}$ and $\mc{F}$ is a set having a single element for feature $x$ which holds a pair 
the first element is $w$ and the second element is a $\mc{Y}$.
%The $\langle C,F\rangle$ creates \emph{two} classes for $x$: one for $x$ itself, the other for feature $x.w$.

%\pn \texttt{Feature} $x$ does not have the same meaning as in (\ref{eq:df-featuredby}).  The differences are that $x$ is a feature of $Z$,  and $x$ is a relation containing a single pair.  
%\comment{(\ref{eq:df-featuremember}) imposes multiplicity/cardinality of one.}

\subsection{Feature with Feature}

\pn Because features are types, they can have their own features.

\begin{restatable}{definition}{dffeatureoffeature}~\blTitle{df-featureoffeature}~   Define \texttt{featured by}
\begin{align}
\texttt{feature}~x~\texttt{typed by}~ Y ~\{~\texttt{feature}~w~\texttt{typed by}~ Z~\}~\Rightarrow  \notag \\
\mc{x}=\{~\langle \textbf{feature},``x",\langle \mc{C},\mc{F}\rangle\rangle~\vert~\mc{C}~=~\{~\langle v,y\rangle~\vert~y \in \mc{Y} \wedge v \in \mc{\mathrm{V}}~\} \notag \\
~\wedge~\mc{F}~=~\{~\langle\textbf{feature},``w",\langle c,z\rangle\rangle~\vert~c~\in~\mc{C}~\wedge~z~\in~\mc{Z}~\}~\}~\label{eq:df-featureoffeature}\tag{df-featureoffeature}
\end{align}
\end{restatable}

\begin{restatable}{leanaxiom}{encodePair}~\leanTitle{encodePair}~ Kuratowski-style pairing, in Lean, encoding
an ordered pair of \texttt{Element}s as a single \texttt{Element}. Needed because features can have features
(this very definition), whose own domain is itself a relation of pairs, and so on to unbounded nesting
depth --- without a uniform way to fold a pair back into a plain \texttt{Element}, each nesting level would
need its own primitive.

\texttt{\textcolor{leancolor}{axiom} \textcolor{leanyellow}{encodePair} : Element $\to$ Element $\to$ Element}
\label{encodePair}
\end{restatable}

\begin{restatable}{leanaxiom}{encodePairInj}~\leanTitle{encodePair\_inj}~ \texttt{encodePair} is injective:
distinct pairs give distinct encodings. No other structure is assumed, matching this project's general
practice of adding only the constraint actually needed.

\texttt{\textcolor{leancolor}{axiom} \textcolor{leanyellow}{encodePair\_inj} \{a b c d : Element\} : encodePair a b = encodePair c d $\to$ a = c $\wedge$ b = d}
\label{encodePair_inj}
\end{restatable}

\begin{restatable}{leanaxiom}{dfFeatureoffeatureLean}~\leanTitle{df\_featureoffeature}~ Define \texttt{featured
by} for nested features in Lean: a domain-less feature $x$ typed by $Y$, with its own owned feature $w$ typed
by $Z$, has content pairing the universe with $Y$'s extension (via \texttt{df\_domainless}), and $w$'s
relation pairs $x$'s own encoded pairs (domain $v$, codomain from $Y$) with $Z$'s extension.

\texttt{\textcolor{leancolor}{axiom} \textcolor{leanyellow}{df\_featureoffeature} \{x w Y' Z' : Element\} \{Cy Cz : Set Element\}}

\texttt{~~(hy : (DesignKind.type, Y', Cy) $\in$ Design) (hz : (DesignKind.type, Z', Cz) $\in$ Design) :}

\texttt{~~DomainlessDecl x Y' $\to$ OwnedFeatureTypedBy x w Z' $\to$}

\texttt{~~~~FeatureContent x \{p : Element $\times$ Element $\vert$ p.2 $\in$ Cy\} $\wedge$}

\texttt{~~~~FeatureMemberOf x w \{q : Element $\times$ Element $\vert$}

\texttt{~~~~~~($\exists$ v y, q.1 = encodePair v y $\wedge$ y $\in$ Cy) $\wedge$ q.2 $\in$ Cz\}}
\label{df_featureoffeature}
\end{restatable}

\pn Feature $\mc{x}$ is domain-less so get default domain of the universe, $\mc{\mathrm{V}}$.  Feature $x$ is still a relation (set of pairs), but has its own member feature $w$, which is also a relation, but its domain is a pair $\langle v,y\rangle$ from $x$ and co-domain $z$.
The $\langle \mc{C},\mc{F}\rangle$ creates \emph{two} classes for $x$: one for $x$ itself, the other for feature $x.w$.
%See (\ref{eq:df-featuremember}) for a simpler type having a feature.

\subsection{Fixed Feature Multiplicity}  

\pn A Feature Member can have multiplicity, in square brackets, \texttt{[ ]}, either following the identifier of the feature, or after the featured type, but not both.
When multiplicity is a natural number, that indicates the cardinality (number of pairs) of the relation.

\ecaption{Fixed Feature Multiplicity}
\begin{lstlisting}
type Person { feature parent[2] typed by Person; }
\end{lstlisting}
says each \bemph{Person} \texttt{Type} must have a relation named \bemph{parent} having exactly two pairs, with the first element being the featuring \bemph{Person}, and the second element being \bemph{different} featured \bemph{Persons}.  

\begin{gather}
\mc{Person} = \{~\langle \textbf{type},w,\mc{C}\rangle~\vert~w=W(``Person")~\wedge \{y_1,y_2 \} \in \mc{Person} \notag \\
~\wedge~ \mc{C}=\{~\langle~\textbf{feature}, ``parent",\mc{R} \rangle~\vert ~\mc{R}=\{~\langle w,y_1\rangle,~\langle w,y_2\rangle~\}~\}~\} \notag
\end{gather}

\pn Default for multiplicity is \texttt{unique}.  (Sets and classes have no repeated elements.  For \texttt{nonunique} something other than sets or classes must be used.)

\begin{restatable}{definition}{dffixedfeaturemultiplicity}~\blTitle{df-fixedfeaturemultiplicity}~   Define \texttt{Fixed Multiplicity}
\begin{align}
\texttt{type}~ Z ~\texttt{\{ feature}~ x[k] ~\texttt{typed by}~ Y~\texttt{; \}}~\Rightarrow  \notag \\
\mc{Z}=\{~\langle \textbf{type},w,\mc{C}\rangle~\vert~w=W(``Z")~\wedge \{y_1,\ldots,y_k \} \in \mc{Y} \notag \\
~\wedge~ \mc{C}=\{~\langle~\textbf{feature}, ``x",\mc{R} \rangle~\vert ~\mc{R}=\{~\langle w,y_1\rangle,\ldots,\langle w,y_k\rangle~\}~\}~\}~\label{eq:df-fixedfeaturemultiplicity}\tag{df-fixedfeaturemultiplicity}
\end{align}
\end{restatable}

\begin{restatable}{leanaxiom}{OwnedFeatureTypedByExact}~\leanTitle{OwnedFeatureTypedByExact}~ \texttt{type Z \{
feature x[k] typed by Y; \}} (exact feature multiplicity), in Lean.

\texttt{\textcolor{leancolor}{axiom} \textcolor{leanyellow}{OwnedFeatureTypedByExact} : Element $\to$ Element $\to$ Element $\to$ $\mathbb{N}$ $\to$ Prop}
\label{OwnedFeatureTypedByExact}
\end{restatable}

\begin{restatable}{leanaxiom}{dfFixedfeaturemultiplicityLean}~\leanTitle{df\_fixedfeaturemultiplicity}~ Define
\texttt{Fixed Multiplicity} in Lean: a feature member declared with exact multiplicity $k$ has a relation of
cardinality exactly $k$, every pair's co-domain drawn from $Y$'s extension. Reuses \texttt{Set.ncard} the same
way \texttt{cardElem}/\texttt{df\_typemultiplicity} did in \S2.2.3, just for pair-sets.

\texttt{\textcolor{leancolor}{axiom} \textcolor{leanyellow}{df\_fixedfeaturemultiplicity} \{Z x Y' : Element\} \{Cy : Set Element\}}

\texttt{~~\{R : Set (Element $\times$ Element)\} \{k : $\mathbb{N}$\} (hy : (DesignKind.type, Y', Cy) $\in$ Design) :}

\texttt{~~OwnedFeatureTypedByExact Z x Y' k $\to$ FeatureMemberOf Z x R $\to$}

\texttt{~~~~R.ncard = k $\wedge$ $\forall$ p $\in$ R, p.2 $\in$ Cy}
\label{df_fixedfeaturemultiplicity}
\end{restatable}

\pn The $y_1,\ldots,y_k$ are \texttt{unique} because $\mc{R}$ is a set and cannot contain duplicate pairs.

\subsection{Unbound Feature Multiplicity}  

\pn Multiplicity of a Feature Member may be unknown, or may be differ between instances, indicated by a star in square brackets, \texttt{[*]}.  Such unknown or undefined multiplicity is \bemph{unbound} not \bemph{infinite}.

\ecaption{Unbound Feature Multiplicity} 
\begin{lstlisting}
type Person { feature children[*] typed by Person; }
\end{lstlisting}
says each Person \texttt{Type} must have a relation named ``children'' having any (natural) number of pairs, with the first element being the featuring Person, and the second element being \bemph{different} featured Persons.  

\begin{restatable}{definition}{dfundefinedfeaturemultiplicity}~\blTitle{df-undefinedfeaturemultiplicity}~ \\ ~~ Define \texttt{Undefined Feature Multiplicity}
\begin{align}
\texttt{type}~ Z ~\texttt{\{ feature}~ x[\star] ~\texttt{typed by}~ Y~\texttt{; \}}~\Rightarrow  \notag \\
\mc{Z}=\{~\langle \textbf{type},w,\mc{C}\rangle~\vert~w=W(``Z")~\wedge~k \in \mathbb{N}_0~\wedge \{ y_1,\ldots,y_k \} \subseteq \mc{Y} \notag \\
~\wedge~ \mc{C}=\{~\langle~\textbf{feature}, ``x",\mc{R} \rangle~\vert ~\mc{R}=\{~\langle w,y_1\rangle,\ldots,\langle w,y_k\rangle~\}~\}~\}~\label{eq:df-undefinedfeaturemultiplicity}\tag{df-undefinedfeaturemultiplicity}
\end{align}
\end{restatable}

\begin{restatable}{leanaxiom}{OwnedFeatureTypedByStar}~\leanTitle{OwnedFeatureTypedByStar}~ \texttt{type Z \{
feature x[*] typed by Y; \}} (unbounded feature multiplicity), in Lean.

\texttt{\textcolor{leancolor}{axiom} \textcolor{leanyellow}{OwnedFeatureTypedByStar} : Element $\to$ Element $\to$ Element $\to$ Prop}
\label{OwnedFeatureTypedByStar}
\end{restatable}

\begin{restatable}{leanaxiom}{dfUndefinedfeaturemultiplicityLean}~\leanTitle{df\_undefinedfeaturemultiplicity}~
Define \texttt{Undefined Feature Multiplicity} in Lean: an unbounded feature member's relation has every
pair's co-domain drawn from $Y$'s extension, with no cardinality constraint at all.

\texttt{\textcolor{leancolor}{axiom} \textcolor{leanyellow}{df\_undefinedfeaturemultiplicity} \{Z x Y' : Element\} \{Cy : Set Element\}}

\texttt{~~\{R : Set (Element $\times$ Element)\} (hy : (DesignKind.type, Y', Cy) $\in$ Design) :}

\texttt{~~OwnedFeatureTypedByStar Z x Y' $\to$ FeatureMemberOf Z x R $\to$ $\forall$ p $\in$ R, p.2 $\in$ Cy}
\label{df_undefinedfeaturemultiplicity}
\end{restatable}

The cardinality of $x$ may be zero or greater.
%\comment{Defining * as `infinity' is erroneous.  No constructed system has infinite anything.}

\subsection{Feature Multiplicity Range}  

\pn Multiplicity may also be defined as a range.  The upper bound of multiplicity range may be undefined.

\ecaption{Feature Multiplicity Range}
\begin{lstlisting}
type Person { feature bff[2..6] typed by Person; }
\end{lstlisting}
says each Person \texttt{Type} must have a relation named ``bff'' having from two to six (inclusive) pairs, with the first element being the featuring Person, and the second element being \bemph{different} featured Persons.  

\begin{restatable}{definition}{dffeaturemultiplicityrange}~\blTitle{df-featuremultiplicityrange}~   Define \texttt{Feature Multiplicity Range}
\begin{align}
\texttt{type}~ Z ~\texttt{\{ feature}~ x[l..u] ~\texttt{typed by}~ Y~\texttt{; \}}~\Rightarrow  \notag \\
\mc{Z}=\{~\langle \textbf{type},w,\mc{C}\rangle~\vert~w=W(``Z")~\wedge~\{ y_l,\ldots,y_u\} \subseteq \mc{Y} \notag \\
~\wedge~ \mc{C}=\{~\langle~\textbf{feature}, ``x",\mc{R} \rangle~\vert ~\mc{R}=\{~\langle w,y_l\rangle,\ldots,\langle w,y_u\rangle~\}~\}~\}\label{eq:df-featuremultiplicityrange}\tag{df-featuremultiplicityrange}
\end{align}
\end{restatable}

\begin{restatable}{leanaxiom}{OwnedFeatureTypedByRange}~\leanTitle{OwnedFeatureTypedByRange}~ \texttt{type Z \{
feature x[l..u] typed by Y; \}} (feature multiplicity range), in Lean.

\texttt{\textcolor{leancolor}{axiom} \textcolor{leanyellow}{OwnedFeatureTypedByRange} : Element $\to$ Element $\to$ Element $\to$ $\mathbb{N}$ $\to$ $\mathbb{N}$ $\to$ Prop}
\label{OwnedFeatureTypedByRange}
\end{restatable}

\begin{restatable}{leanaxiom}{dfFeaturemultiplicityrangeLean}~\leanTitle{df\_featuremultiplicityrange}~ Define
\texttt{Feature Multiplicity Range} in Lean: a range-bounded feature member's relation has cardinality
between $l$ and $u$, every pair's co-domain drawn from $Y$'s extension.

\texttt{\textcolor{leancolor}{axiom} \textcolor{leanyellow}{df\_featuremultiplicityrange} \{Z x Y' : Element\} \{Cy : Set Element\}}

\texttt{~~\{R : Set (Element $\times$ Element)\} \{l u : $\mathbb{N}$\} (hy : (DesignKind.type, Y', Cy) $\in$ Design) :}

\texttt{~~OwnedFeatureTypedByRange Z x Y' l u $\to$ FeatureMemberOf Z x R $\to$}

\texttt{~~~~l $\le$ R.ncard $\wedge$ R.ncard $\le$ u $\wedge$ $\forall$ p $\in$ R, p.2 $\in$ Cy}
\label{df_featuremultiplicityrange}
\end{restatable}

\pn The upper bound of multiplicity range may also be undefined.

\begin{restatable}{definition}{dffeaturemultiplicityrangestar}~\blTitle{df-featuremultiplicityrangestar}~ \\~~  Define \texttt{Feature Multiplicity Range Star}
\begin{align}
\texttt{type}~ Z ~\texttt{\{ feature}~ x[l..\star] ~\texttt{typed by}~ Y~\texttt{; \}}~\Rightarrow  \notag \\
\mc{Z}=\{~\langle \textbf{type},w,\mc{C}\rangle~\vert~w=W(``Z")~\wedge~k \in \mathbb{N}_0~\wedge \{ y_l,\ldots,y_k \} \subseteq \mc{Y} \notag \\
~\wedge~ \mc{C}=\{~\langle~\textbf{feature}, ``x",\mc{R} \rangle~\vert ~\mc{R}=\{~\langle w,y_l\rangle,\ldots,\langle w,y_k\rangle~\}~\}~\}~~\label{eq:df-featuremultiplicityrangestar}\tag{df-featuremultiplicityrangestar}
\end{align}
\end{restatable}

\begin{restatable}{leanaxiom}{OwnedFeatureTypedByRangeStar}~\leanTitle{OwnedFeatureTypedByRangeStar}~
\texttt{type Z \{ feature x[l..*] typed by Y; \}} (feature multiplicity range, unbounded upper end), in Lean.

\texttt{\textcolor{leancolor}{axiom} \textcolor{leanyellow}{OwnedFeatureTypedByRangeStar} : Element $\to$ Element $\to$ Element $\to$ $\mathbb{N}$ $\to$ Prop}
\label{OwnedFeatureTypedByRangeStar}
\end{restatable}

\begin{restatable}{leanaxiom}{dfFeaturemultiplicityrangestarLean}~\leanTitle{df\_featuremultiplicityrangestar}~
Define \texttt{Feature Multiplicity Range Star} in Lean: same as \texttt{df\_featuremultiplicityrange} above
but with only a lower bound $l$.

\texttt{\textcolor{leancolor}{axiom} \textcolor{leanyellow}{df\_featuremultiplicityrangestar} \{Z x Y' : Element\} \{Cy : Set Element\}}

\texttt{~~\{R : Set (Element $\times$ Element)\} \{l : $\mathbb{N}$\} (hy : (DesignKind.type, Y', Cy) $\in$ Design) :}

\texttt{~~OwnedFeatureTypedByRangeStar Z x Y' l $\to$ FeatureMemberOf Z x R $\to$}

\texttt{~~~~l $\le$ R.ncard $\wedge$ $\forall$ p $\in$ R, p.2 $\in$ Cy}
\label{df_featuremultiplicityrangestar}
\end{restatable}

\subsection{Type Membership Feature}
\comment{Feature Member and Membership Feature are easily confused.  I'm not sure how to represent \texttt{owned} or not in ZFC.}

\pn If a feature declaration in the body of type is preceded by the keyword \texttt{member}, then the feature is owned by the containing type via a membership relationship, not a feature membership. In this case, the feature is not an owned feature of the containing type, and it does not automatically have the containing type as a featuring type, though it may have featuring types declared in its \texttt{featured by} list 
(see KerML \S7.3.4.1 on declaring the owned typings of a feature).

\ecaption{Member Feature}
\begin{lstlisting}
type A; type C;
type B {
  // Feature f has B as its featuring type, and C as its featured type. 
  feature f typed by C;
  // Feature g has A as its featuring type, not B.
  member feature g typed by C featured by A; }
\end{lstlisting}

\begin{restatable}{definition}{dfmembershipfeature}~\blTitle{df-membershipfeature}~   Define Membership.
\begin{align}
\texttt{type}~ B~\{ \texttt{member feature}~ g ~\texttt{typed by}~ C~\texttt{featured by}~A~\texttt{; \}}~\Rightarrow  \notag \\
\mc{B}=\{~\langle \textbf{type},w,\mc{M}\rangle~\vert~w=W(``B")~\wedge c \in \mc{C}~\wedge a \in \mc{A} \notag \\
~\wedge~ \mc{M}=\{~\langle~\textbf{member}, ``g",\mc{R} \rangle~\vert ~\mc{R}=\{~\langle a,c\rangle~\}~\}~\}
\label{eq:df-membershipfeature}\tag{df-membershipfeature}
\end{align}
\end{restatable}

\begin{restatable}{leanaxiom}{MemberOf}~\leanTitle{MemberOf}~ \texttt{type B \{ member feature g typed by C
featured by A; \}} (\texttt{member feature}, not an owned feature member), in Lean. A fresh
\texttt{Element}-level primitive: $B$ has ordinary (non-feature) membership of $g$, distinct from
\texttt{FeatureMemberOf}/\texttt{OwnedFeatureTypedBy} above --- matching the book's own point that $g$ is
\emph{not} an owned feature of $B$.

\texttt{\textcolor{leancolor}{axiom} \textcolor{leanyellow}{MemberOf} : Element $\to$ Element $\to$ Prop}
\label{MemberOf}
\end{restatable}

\begin{restatable}{leantheorem}{dfMembershipfeatureLean}~\leanTitle{df\_membershipfeature}~ Define Membership
in Lean. A member feature $g$'s own relation is determined entirely by its own \texttt{featured by}
declaration, independent of its containing type $B$ --- exactly \texttt{df\_featuredby}'s own conclusion,
applied to $g$. \texttt{MemberOf B g} is included as a hypothesis only for faithfulness to the concrete-syntax
shape ($B$ does own $g$ as a \emph{member}), not because the proof needs it: that $B$'s membership plays no
role in $g$'s content is precisely the book's own point (``the feature is not an owned feature of the
containing type''), so this one genuinely is a theorem.

\texttt{\textcolor{leancolor}{theorem} \textcolor{leanyellow}{df\_membershipfeature} \{B g C' A' : Element\} \{Cc Ca : Set Element\}}

\texttt{~~(hc : (DesignKind.type, C', Cc) $\in$ Design) (ha : (DesignKind.type, A', Ca) $\in$ Design)}

\texttt{~~(\_hmem : MemberOf B g) (hfbd : FeaturedByDecl g C' A') :}

\texttt{~~FeatureContent g \{p : Element $\times$ Element $\vert$ p.1 $\in$ Ca $\wedge$ p.2 $\in$ Cc\} :=}

\texttt{~~df\_featuredby hc ha hfbd}
\label{df_membershipfeature}
\end{restatable}

\subsection{Multiple Type Body Elements}

\pn A \texttt{TypeBody} may have multiple elements of various kinds.  Each element is semantically represented by a triple, the first element telling what kind of element it is, the second element containing a string with its identifier, and the third element it ZFC-class definition.
However, rather than the type's (single) ZFC-class definition, it has a tuple with an entry for each type body element.

\pn In the definition, \texttt{feature} elements are used, without loss of generality (substituting the appropriate triple).

\ecaption{Type with Multiple Features}
\begin{lstlisting}
type Z {
  feature x1 typed by Y1;
  feature x2 typed by Y2;
  . . .
  feature xj typed by Yj;
}
\end{lstlisting}

\begin{restatable}{definition}{dftypebody}~\blTitle{df-typebody}~   Define multiple \texttt{TypeBody} elements.
\begin{align}
\texttt{type}~ Z~\{~ \texttt{feature}~ x_1 ~\texttt{typed by}~ Y_1~\texttt{;}~\ldots~
  \texttt{feature}~ x_j ~\texttt{typed by}~ Y_j~\texttt{;}\}~\Rightarrow  \notag \\
\mc{Z}=\{~\langle \textbf{type},w,\langle \mc{C_1},\ldots,\mc{C_j}\rangle~\rangle~\vert~w=W(``Z")~\wedge y_1 \in \mc{Y_1} \wedge\ldots\wedge y_j \in \mc{Y_j} \notag \\
~\wedge~ \mc{C_1}=\{~\langle~\textbf{feature}, ``x_1",\mc{R_1} \rangle~\vert ~\mc{R_1}=\{~\langle w,y_1\rangle~\}~\} \notag \\
~~~\wedge~\ldots~\wedge %\notag \\
~~\mc{C_j}=\{~\langle~\textbf{feature}, ``x_j",\mc{R_j} \rangle~\vert ~\mc{R_j}=\{~\langle w,y_j\rangle~\}~\} ~\}
\label{eq:df-typebody}\tag{df-typebody}
\end{align}
\end{restatable}

\begin{restatable}{leanaxiom}{OwnedFeatures}~\leanTitle{OwnedFeatures}~ \texttt{type Z \{ feature x1 typed by
Y1; ...; feature xj typed by Yj; \}} (\texttt{TypeBody} with several feature members, variable arity $j$), in
Lean. \texttt{feats} holds the declared (feature, type) pairs as a \texttt{List}, sidestepping Lean's lack of
a literal ``...'' ellipsis for a fixed but unspecified arity $j$.

\texttt{\textcolor{leancolor}{axiom} \textcolor{leanyellow}{OwnedFeatures} : Element $\to$ List (Element $\times$ Element) $\to$ Prop}
\label{OwnedFeatures}
\end{restatable}

\begin{restatable}{leanaxiom}{dfTypebodyLean}~\leanTitle{df\_typebody}~ Define multiple \texttt{TypeBody}
elements in Lean: generalizes \texttt{df\_featuremember} (domain $Z$'s own extension, since no
\texttt{specializes} clause is given here) to every feature $(x_i,Y_i)$ declared in \texttt{feats}
independently --- the ``...'' in the book's own $C_1,\ldots,C_j$ really means ``this holds
member-by-member'', which is exactly what quantifying over \texttt{feats} states.

\texttt{\textcolor{leancolor}{axiom} \textcolor{leanyellow}{df\_typebody} \{Z : Element\} \{feats : List (Element $\times$ Element)\} \{Cz : Set Element\}}

\texttt{~~(hz : (DesignKind.type, Z, Cz) $\in$ Design) :}

\texttt{~~OwnedFeatures Z feats $\to$}

\texttt{~~~~$\forall$ p $\in$ feats, $\exists$ Cy : Set Element, (DesignKind.type, p.2, Cy) $\in$ Design $\wedge$}

\texttt{~~~~~~FeatureMemberOf Z p.1 \{q : Element $\times$ Element $\vert$ q.1 $\in$ Cz $\wedge$ q.2 $\in$ Cy\}}
\label{df_typebody}
\end{restatable}

\begin{description}
\item[Unique: ] Sets and classes disallow repeated elements, so unique is appropriate default for features.
\item[Non-unique: ] A feature may be \texttt{nonunique} then for any \bemph{domain} instance, the same \bemph{co-domain} instance may appear more than once as a value of the feature.
Non-unique beaks the definition of features as relations, and is nonsensical except in terms of a stream of values over time, such as sensor measurements, which will be defined temporally. 
\item[Ordered: ] A feature is ordered if its \bemph{co-domain} is ordered.
\end{description}

\begin{restatable}{definition}{dfordered}~\blTitle{df-ordered}~  Define \texttt{Ordered}
\begin{align}
\texttt{type}~ Z ~\texttt{\{ feature}~ x[\star] ~\texttt{ordered typed by}~ Y~\texttt{; \}}~\Rightarrow  \notag \\
\forall \langle z,y_1\rangle, \langle z,y_2\rangle~\in~Z.x~\vert~(y_1<y_2)~\rightarrow ( \langle z,y_1\rangle< \langle z,y_2\rangle)~\label{eq:df-ordered}\tag{df-ordered}
\end{align}
\end{restatable}

\begin{restatable}{leanaxiom}{OrderedOn}~\leanTitle{OrderedOn}~ \texttt{OrderedOn Y a b} means $a<b$ under
$Y$'s own (type-relative) ordering, in Lean --- not every type has a natural order (the book's own note: ``if
$Y$ is not ordered, the feature $x$ cannot be ordered''), so this is parameterized by $Y$ rather than a single
global order on \texttt{Element}.

\texttt{\textcolor{leancolor}{axiom} \textcolor{leanyellow}{OrderedOn} : Element $\to$ Element $\to$ Element $\to$ Prop}
\label{OrderedOn}
\end{restatable}

\begin{restatable}{leanaxiom}{OrderedDecl}~\leanTitle{OrderedDecl}~ \texttt{type Z \{ feature x[*] ordered
typed by Y; \}}, in Lean.

\texttt{\textcolor{leancolor}{axiom} \textcolor{leanyellow}{OrderedDecl} : Element $\to$ Element $\to$ Element $\to$ Prop}
\label{OrderedDecl}
\end{restatable}

\begin{restatable}{leandefinition}{PairOrderedOn}~\leanTitle{PairOrderedOn}~ \texttt{df-ordered}'s pair order,
in Lean: pairs sharing the same first component are ordered by their second component's own ($Y$-relative)
order. A genuine \texttt{def}, not an axiom --- this is exactly what ``the pair inherits $y_1$'s/$y_2$'s order
when $z$ is fixed'' means, not an independent constraint.

\texttt{\textcolor{leancolor}{def} \textcolor{leanyellow}{PairOrderedOn} (Y : Element) (p q : Element $\times$ Element) : Prop := p.1 = q.1 $\wedge$ OrderedOn Y p.2 q.2}
\label{PairOrderedOn}
\end{restatable}

\begin{restatable}{leantheorem}{dfOrderedLean}~\leanTitle{df\_ordered}~ Define \texttt{Ordered} in Lean: when
two pairs of an ordered feature's relation share the same first component $z$, and their second components
are $Y$-ordered, the pairs themselves are $Y$-ordered (via \texttt{PairOrderedOn}). Free directly from
\texttt{PairOrderedOn}'s own definition --- genuinely a theorem, matching the
\texttt{df\_typemultiplicity}-family precedent of only axiomatizing what's not already definitional.

\texttt{\textcolor{leancolor}{theorem} \textcolor{leanyellow}{df\_ordered} \{Z x Y' : Element\} \{R : Set (Element $\times$ Element)\} \{z y1 y2 : Element\}}

\texttt{~~(\_hR : FeatureMemberOf Z x R) (\_hDecl : OrderedDecl Z x Y')}

\texttt{~~(\_h1 : (z, y1) $\in$ R) (\_h2 : (z, y2) $\in$ R) (hlt : OrderedOn Y' y1 y2) :}

\texttt{~~PairOrderedOn Y' (z, y1) (z, y2) := $\langle$rfl, hlt$\rangle$}
\label{df_ordered}
\end{restatable}

Of course, if the co-domain, $Y$ is not ordered, the feature $x$ cannot be ordered.  The relation $<$ is a total ordering on $Y$ iff:
\begin{itemize}
\item $\forall y_1, y_2 \in Y~\vert~y_1\ne y_2~\rightarrow y_1<y_2~\vee~y_2<y_1$
\item $\forall y_1, y_2, y_3 \in Y~\vert~y_1<y_2~\wedge~y_2<y_3~\rightarrow~y_1<y_3$
\end{itemize}

\subsection{Feature Direction and Conjugation}\label{sec:direction}

\pn Feature \bemph{direction} \texttt{in}, \texttt{out}, and \texttt{inout} determines what is allowed to change values of a feature's co-domain.  This applies to instances of a type having a directed feature.  Conjugation of a type, \texttt{\midtilde}, reverses direction of its features.
Semantics of directed features requires different values at different times is defined in terms of modal logic in section \ref{sec:variables}.

\subsection{Type Specialization with Features} 

\pn  The class definition of a type $Z$ having feature $x$ requires elements of $Z$ to have a relation $Z.x$, but may have other things as well.
(\ref{eq:df-specializes}) defines specialization as subclassing.  Consider some other type $Y$, having a feature $x$ defined the same way as $Z$, but also having a feature $u$ (defined its own way) will be a subclass of $Z$:  $Y~\subset~Z$.

\pn Let type $Z$ have a feature $x$ typed by $A$, and type $Y$ which specializes $Z$ have a feature $u$ typed by $B$:

\ecaption{Type Specialization with Features}
\begin{lstlisting}
type Z { feature x  typed by A ; }
type Y specializes Z { feature u typed by B ; }
\end{lstlisting}

then the class definition of $Y$ will be:
\begin{restatable}{definition}{dfspecializationwithfeatures}~\blTitle{df-specializationwithfeatures}~  \\ ~~ Define \texttt{Type Specialization with Features}
\begin{align}
\texttt{type}~ Z ~\texttt{\{ feature}~ x~\texttt{typed by}~A~\texttt{; \}} \notag \\ %\notag \\
\texttt{type}~ Y ~\texttt{specializes}~Z~\texttt{\{ feature}~ u~\texttt{typed by}~B~\texttt{; \}} ~\Rightarrow  \notag \\
\mc{Y}=\{~\langle \textbf{type},w,\langle \mc{C_z},\mc{C_u},\mc{C_x}\rangle~\rangle~\vert~w=W(``Y")~\notag \\
\wedge ~a \in \mc{A} ~\wedge~ b \in \mc{B}
~\wedge~\mc{C_z}=\{~z\in \mc{Z}~\} \notag \\
~\wedge~ \mc{C_u}=\{~\langle~\textbf{feature}, ``u",\mc{R_u} \rangle~\vert ~\mc{R_u}=\{~\langle w,b\rangle~\}~\} \notag \\
~~~\wedge~\mc{C_x}=\{~\langle~\textbf{feature}, ``x",\mc{R_x} \rangle~\vert ~\mc{R_x}=\{~\langle w,a\rangle~\}~\} ~\}
\label{eq:df-specializationwithfeatures}\tag{df-specializationwithfeatures}
\end{align}
\end{restatable}

\begin{restatable}{leantheorem}{dfSpecializationwithfeaturesLean}~\leanTitle{df\_specializationwithfeatures}~
Define \texttt{Type Specialization with Features} in Lean: a type $Y$ specializing $Z$, with its own owned
feature $u$ typed by $B$, has a feature member for $u$ whose relation pairs $Z$'s extension (domain, since
$Y$ specializes $Z$) with $B$'s extension --- exactly \texttt{df\_featuremember} applied with $T:=Z$, so this
one genuinely is a theorem. Scoped down from the book's own full $\langle C_z,C_u,C_x\rangle$ triple: the
re-inclusion of $Z$'s own inherited feature $x$ as part of $Y$'s content isn't covered here, since nothing in
this project axiomatizes feature inheritance through \texttt{Specializes} itself.

\texttt{\textcolor{leancolor}{theorem} \textcolor{leanyellow}{df\_specializationwithfeatures} \{Z Y u B : Element\} \{Cz Cb : Set Element\}}

\texttt{~~(hz : (DesignKind.type, Z, Cz) $\in$ Design) (hb : (DesignKind.type, B, Cb) $\in$ Design)}

\texttt{~~(hspec : Specializes Y Z) (hyu : OwnedFeatureTypedBy Y u B) :}

\texttt{~~FeatureMemberOf Y u \{p : Element $\times$ Element $\vert$ p.1 $\in$ Cz $\wedge$ p.2 $\in$ Cb\} :=}

\texttt{~~df\_featuremember hz hb hspec hyu}
\label{df_specializationwithfeatures}
\end{restatable}

\pn Therefore $\mc{Y}~\subset~\mc{Z}$ because the rule defining class $\mc{Y}$, implies the rule defining class $\mc{Z}$.
The rule defining the ZFC-class for $\mc{Z}$ will have the terms for $w$, $a$, and $\mc{C_x}$, which is implied by the rule for $Y$ in (\ref{eq:df-specializationwithfeatures}) by (\ref{eq:simpl}).  \comment{Full exposition and proof left to the reader.}

\pn Specialization is generalized for any number of features, and specialization of specialization the same way.

\subsection{Feature Inverting}  

\pn Feature inverting exchanges domain with co-domain of the relation of a specific feature.
Feature inverting does \emph{not} invert the relation itself.

\pn Where Y and Z are names of \texttt{Types}:
\ecaption{Feature Inverting} 
\begin{lstlisting}
feature x typed by Y featured by Z;
inverse w of x;
\end{lstlisting}
says $\mc{x}$ is a class of pairs where the first element is from $\mc{Z}$ and the second element is from $\mc{Y}$, and $\mc{w}$ is a class of pairs where the first element is from $\mc{Y}$ and the second element is from $\mc{Z}$.


\begin{restatable}{definition}{df-featureinverting}~\blTitle{df-featureinverting}~   Define \texttt{Feature Inverting}
\begin{align}
\texttt{feature}~x~\texttt{typed by}~ Y ~\texttt{featured by}~ Z~\texttt{;}~\notag \\ 
\texttt{inverse}~ w ~\texttt{of} ~x~\texttt{;}~\Rightarrow  \notag \\
~\mc{x} \subseteq~\{~\langle z,y\rangle~\vert~y \in \mc{Y} \wedge z \in \mc{Z}~\}~\wedge~w \subseteq~\{~\langle y,z\rangle~\vert~y \in \mc{Y} \wedge z \in \mc{Z}~\}
\label{eq:df-featureinverting}\tag{df-featureinverting}
\end{align}
\end{restatable}

\begin{restatable}{leantheorem}{dfFeatureinvertingLean}~\leanTitle{df\_featureinverting}~ Define \texttt{Feature
Inverting} in Lean: ``\texttt{inverse w of x}'' needs no new primitive at all --- it's captured directly as
$w$ independently satisfying \texttt{FeaturedByDecl} with domain/co-domain swapped relative to $x$, so both
halves reduce straight to \texttt{df\_featuredby}, applied twice. The book's own formula states $\subseteq$,
but \texttt{df\_featuredby} already gives the stronger exact content --- that stronger fact is given directly
here rather than artificially weakening it, so this one genuinely is a theorem.

\texttt{\textcolor{leancolor}{theorem} \textcolor{leanyellow}{df\_featureinverting} \{x w y z : Element\} \{Cy Cz : Set Element\}}

\texttt{~~(hy : (DesignKind.type, y, Cy) $\in$ Design) (hz : (DesignKind.type, z, Cz) $\in$ Design)}

\texttt{~~(hxfb : FeaturedByDecl x y z) (hwfb : FeaturedByDecl w z y) :}

\texttt{~~FeatureContent x \{p : Element $\times$ Element $\vert$ p.1 $\in$ Cz $\wedge$ p.2 $\in$ Cy\} $\wedge$}

\texttt{~~FeatureContent w \{p : Element $\times$ Element $\vert$ p.1 $\in$ Cy $\wedge$ p.2 $\in$ Cz\} :=}

\texttt{~~$\langle$df\_featuredby hy hz hxfb, df\_featuredby hz hy hwfb$\rangle$}
\label{df_featureinverting}
\end{restatable}

\subsection{Owned Feature Inverting}  

\pn Inverting an owned feature is trickier than general feature inverting.  Consider an example from KerML \S7.3.4.7:
\ecaption{Owned Feature Inverting} 
\begin{lstlisting}
type Person {
  feature children typed by Person[*];
  feature parents typed by Person[*] inverse of children;
  }
\end{lstlisting}
\pn The intent would seem to be that if Gail was a child of Paul, then Paul must be the parent of Gail.  However, \ref{eq:df-featureinverting}
says that parents merely exchanges domain and co-domain, both of which are Person.  It does not invert the relation.  Perhaps instead we want:
\begin{lstlisting}
type Person; 
feature children typed by Person featured by Person;
parents inverse of children;
\end{lstlisting}





%\pn Given the definition of feature multiplicity ~(\ref{eq:df-undefinedfeaturemultiplicity}), \texttt{parents} cannot be the inverse of \texttt{children}!  
%\comment{This is so obvious, does it really need to be explained?}  
%
%\pn Only owned features with multiplicity of one can be inverted:
%\begin{lstlisting}
%type Person {
%  feature child typed by Person;
%  feature parent typed by Person inverse of child;
%  }
%\end{lstlisting}
%\pn Now \texttt{parent} is sensibly the inverse of \texttt{child} by swapping domain and co-domain of ~\ref{eq:df-featureinverting}.

\subsection{Feature of Feature}  

\pn Although it is confusing, features themselves may have features.  Features are types after all.  Nested features have qualified names reflecting their nesting.
Consider:
\ecaption{Feature of Feature} 
\begin{lstlisting}
type A { 
  feature b1: B {
    feature c1: C; 
  }
}
\end{lstlisting}
%type C {
%  feature b2: B {
%    feature a2: A inverse of A::b1::c1;
%  } 
%}

\pn The domain of a nested \texttt{Feature} is its owning \texttt{Feature}.

\pn By \ref{eq:df-featuremember}, feature \texttt{b1} is a set of pairs, $\{~\langle a,b\rangle~\}$ such that $a \in A$ and $b \in B$.
Also by \ref{eq:df-featuremember} feature \texttt{c1} is a set of pairs $\{~\langle \langle a,b \rangle, c\rangle~\}$ such that $\langle a,b\rangle\in b_1$ and $c\in C$.
The domain of \texttt{c1} is a set of pairs--not $A$ as one might expect.

\comment{Make formal definition?}

%\comment{This is stated nowhere.  The abstract syntax for \texttt{Feature} extends \texttt{Type}, so the domain of \texttt{b1} should be \texttt{B}--not \texttt{A}.
%I will not bother to formally define this, because the statement above is clear enough.}

%Feature \texttt{a2} is defined using the qualified name of feature \texttt{c1}.  However, feature \texttt{c1} is a relation consisting of a single pair with domain \texttt{B} and co-domain \texttt{C}, and feature \texttt{a2} is a relation consisting of a single pair with domain \texttt{B} and co-domain \texttt{A}.  So \texttt{a2} \bemph{cannot} be the inverse of \texttt{c1}!

\subsection{Feature Chaining} 

\pn  Features can be chained by separating their identifiers with periods.  Feature chaining is composition of their relations.

\begin{restatable}{definition}{dffeaturechaining}~\blTitle{df-featurechaining}~   Define \texttt{Feature Chaining}\footnote{The ordering of the relation composition operator $\circ$ is the reverse of what you might intuit.}
\begin{align}
\texttt{feature}~a~\texttt{chains}~ b.c ~\texttt{;}~\Rightarrow  \notag \\
\vdash a~=~(c~\circ~b)
\label{eq:df-featurechaining}\tag{df-featurechaining}
\end{align}
\end{restatable}

\begin{restatable}{leanaxiom}{ChainsDecl}~\leanTitle{ChainsDecl}~ \texttt{feature a chains b.c;}, in Lean. A
fresh \texttt{Element}-level primitive, same reasoning as \texttt{Specializes}.

\texttt{\textcolor{leancolor}{axiom} \textcolor{leanyellow}{ChainsDecl} : Element $\to$ Element $\to$ Element $\to$ Prop}
\label{ChainsDecl}
\end{restatable}

\begin{restatable}{leanaxiom}{dfFeaturechainingLean}~\leanTitle{df\_featurechaining}~ Define \texttt{Feature
Chaining} in Lean: a feature $a$ declared as chaining $b.c$ has content the composition of $c$'s and $b$'s own
relations (via \texttt{relCompose}, defined just below at \texttt{df-co}).

\texttt{\textcolor{leancolor}{axiom} \textcolor{leanyellow}{df\_featurechaining} \{a b c : Element\} \{Rb Rc : Set (Element $\times$ Element)\}}

\texttt{~~(hb : FeatureContent b Rb) (hc : FeatureContent c Rc) :}

\texttt{~~ChainsDecl a b c $\to$ FeatureContent a (relCompose Rc Rb)}
\label{df_featurechaining}
\end{restatable}

\begin{restatable}{definition}{dfco}~\blTitle{df-co}~   Define composition of two relations.\footnote{The ordering of the relation composition operator $\circ$ is the reverse of what you might intuit.}
\begin{align}
\vdash (\mc{A}~\circ~\mc{B})~\equiv~\{~\langle x,y\rangle~\vert~\exists z ( x \mc{B} z~\wedge~z \mc{A} y )~\}
\label{eq:df-co}\tag{df-co}
\end{align}
\end{restatable}

\begin{restatable}{leandefinition}{relCompose}~\leanTitle{relCompose}~ Define composition of two relations in
Lean. A genuine \texttt{def}, no axiom needed --- pure set-builder over already-available \texttt{Set}/pair
machinery.

\texttt{\textcolor{leancolor}{def} \textcolor{leanyellow}{relCompose} (A B : Set (Element $\times$ Element)) : Set (Element $\times$ Element) :=}

\texttt{~~\{p : Element $\times$ Element $\vert$ $\exists$ z, (p.1, z) $\in$ B $\wedge$ (z, p.2) $\in$ A\}}
\label{relCompose}
\end{restatable}

\pn When chaining \texttt{b.c} the co-domain of \texttt{b} must be a subset of the domain of \texttt{c} or the chain will be empty.

\pn There is no limit to feature chaining.  Therefore (\ref{eq:df-featurechaining}) may need to be applied successively.

%\pn \paragraph{Composite and Portion:} Features may be declared \texttt{composite} or \texttt{portion} to indicate its existence ceases when its owner ceases to exist.  Therefore they should not be used with Types (which are immutable over time), but Classifiers which have temporal semantics.
%Because these are temporal, definition will be deferred until \ref{par:cnp}.
%
%\pn \paragraph{Derived:}  Features defined as \texttt{derived} have their values determined by Relationships with other Entities.
%
%\pn \paragraph{Readonly:}  Features defined as \texttt{readonly} are constant, which is default.  
%\comment{Should \texttt{readonly} be eliminated?}
%
%\pn \paragraph{In, Out, Inout, and Var:}  These prefixes identify features whose values (elements of the co-domain of the relation) may change over time.  Therefore they should not be used with Types (which are immutable over time), but Classifiers which have temporal semantics.
%Because these are temporal, definition will be deferred until \ref{par:variables}.




